\makeabstract
{
Analyzing the functional role of Adad1
}
{
Karol Diaz (1,2), Kavita Venkataramani (1), Kellee Siegfried (1)
}
{
\textsuperscript{1}Biology Department, University of Massachusetts Boston, Boston, MA\\
\textsuperscript{2}Biology Department, Amherst College, Amherst, MA
}
{
There is an incomplete understanding of how genetic factors lead to infertility. We use zebrafish as a model organism to learn more about these factors. Adad1 (Adenosine deaminase domain-containing protein 1) is a gene that has been shown to affect male fertility, and previous research in our lab has found that Adad1 is essential for germ cell maintenance in zebrafish, with a null mutant leading to total loss of germ cells. This gene encodes a double-stranded RNA-binding motif (dsRBM) and an adenosine deaminase domain (DD). We hypothesize that the dsRBM and DD domain each play a significant role in the functional activity of Adad1. Ultimately, we hope to gain insight into RNA-binding proteins’ role in germ cell maintenance in both zebrafish and humans. 
}
{
We performed an in vitro assay in zebrafish embryos, looking at the RNA regulation by Adad1 and the effects of individual domain knockouts on protein function. In this tethering assay, we tethered our protein of interest to MS2, a bacteriophage coat protein that binds to RNA stem loop helices. When mixed with firefly Luciferase mRNA with these motifs, it leads to a measurable luminescence output after allowing translation to occur. We synthesized the RNA and injected it into zebrafish embryos at the one-cell stage. The tethering assay was repeated with Adad1 domain knockouts. 
}
{
Our initial tethering assay revealed that Adad1 had a significantly higher luciferase protein produced than our control (p <0.05). Based on our hypothesis, we expected that if one of the domains were necessary for Adad1’s functional activity, it would produce lower luciferase proteins. The tethering assay with the Adad1 knockouts significantly differed from MS2, our baseline luminescence, except for the effector protein that only had Adad1 DD. 
}
{
Neither domain seems necessary for the effect Adad1 has on proteins. Our preliminary results indicate that Adad1 interacts with RNA to increase protein production, but future experiments are needed to determine the exact mechanism. 
}

\newpage
